\documentclass[a4paper,10pt,twoside]{article}

\usepackage{a4wide}
\usepackage[utf8]{inputenc}
\usepackage[spanish,es-ucroman,es-noquoting]{babel}
\usepackage{setspace}
\usepackage{fancyhdr}
\usepackage{lastpage}
\usepackage{amsmath}
\usepackage{amsfonts}
\usepackage{verbatim}
\usepackage{graphicx}
\usepackage{float}
\usepackage{algorithmic}
\usepackage{tikz}
\usetikzlibrary{calc}
\usetikzlibrary{decorations.pathreplacing}


% Evita que el documento se estire verticalmente para ocupar
% el espacio vacío en cada página.
\raggedbottom


%%%%%%%%%% Configuración de Fancyhdr - Inicio %%%%%%%%%%
\pagestyle{fancy}
\thispagestyle{fancy}
\lhead{Trabajo Práctico 2, Organización del Computador II}
\rhead{Belloli, Lovisolo, Petaccio}
\renewcommand{\footrulewidth}{0.4pt}
\cfoot{\thepage /\pageref{LastPage}}

\fancypagestyle{caratula} {
   \fancyhf{}
   \cfoot{\thepage /\pageref{LastPage}}
   \renewcommand{\headrulewidth}{0pt}
   \renewcommand{\footrulewidth}{0pt}
}
%%%%%%%%%% Configuración de Fancyhdr - Fin %%%%%%%%%%


%%%%%%%%%% Configuración de Algorithmic - Inicio %%%%%%%%%%
% Entorno propio para customizar la presentación del pseudocódigo
\newenvironment{pseudocodigo}
    {\vspace{0.5em} \begin{algorithmic}}
    {\end{algorithmic} \vspace{0.5em}}

% Alinear comentarios a la derecha
\renewcommand{\algorithmiccomment}[1]{\hfill \{#1\}}
%%%%%%%%%% Configuración de Algorithmic - Fin %%%%%%%%%%


%%%%%%%%%% Macros de tikz - Inicio %%%%%%%%%%
% Uso: \registroCuatro{etiqueta}{x}{y}{a4}{a3}{a2}{a1}
\newcommand{\registroCuatro}[7]{
    \ifthenelse{\equal{#1}{}}{}{
        \draw (#2, {#3 + 0.5}) node[anchor=east]{#1};
    }

    \draw   (#2, #3) rectangle +(4, 1) +(2, 0.5) node{#4}
          ++(4, 0)   rectangle +(4, 1) +(2, 0.5) node{#5}
          ++(4, 0)   rectangle +(4, 1) +(2, 0.5) node{#6}
          ++(4, 0)   rectangle +(4, 1) +(2, 0.5) node{#7};          
}

\newcommand{\registroOcho}[9]{
    \def\etiqueta{#1}
    \def\x{#2}
    \def\y{#3}
    \def\aviii{#4}
    \def\avii{#5}
    \def\avi{#6}
    \def\av{#7}
    \def\aiv{#8}
    \def\aiii{#9}
    \registroOchoX    
}
\newcommand{\registroOchoX}[2]{
    \def\aii{#1}
    \def\ai{#2}
    \ifthenelse{\equal{\etiqueta}{}}{}{
        \draw (\x, {\y + 0.5}) node[anchor=east]{\etiqueta};
    }
    \filldraw[fill=white]
        (\x, \y) rectangle +(2, 1) +(1, 0.5) node{\aviii}
        ++(2, 0) rectangle +(2, 1) +(1, 0.5) node{\avii}
        ++(2, 0) rectangle +(2, 1) +(1, 0.5) node{\avi}
        ++(2, 0) rectangle +(2, 1) +(1, 0.5) node{\av}
        ++(2, 0) rectangle +(2, 1) +(1, 0.5) node{\aiv}
        ++(2, 0) rectangle +(2, 1) +(1, 0.5) node{\aiii}
        ++(2, 0) rectangle +(2, 1) +(1, 0.5) node{\aii}
        ++(2, 0) rectangle +(2, 1) +(1, 0.5) node{\ai};
}


% Uso: \registroDieciseis{etiqueta}{x}{y}{a16}{a15}{a14}...{a1}
\newcommand{\registroDieciseis}[9]{
    \def\etiqueta{#1}
    \def\x{#2}
    \def\y{#3}
    \def\axvi{#4}
    \def\axv{#5}
    \def\axiv{#6}
    \def\axiii{#7}
    \def\axii{#8}
    \def\axi{#9}
    \registroDieciseisX
}
\newcommand{\registroDieciseisX}[9]{ % Auxiliar - no usar directamente
    \def\ax{#1}
    \def\aix{#2}
    \def\aviii{#3}
    \def\avii{#4}
    \def\avi{#5}
    \def\av{#6}
    \def\aiv{#7}
    \def\aiii{#8}
    \def\aii{#9}
    \registroDieciseisXX
}
\newcommand{\registroDieciseisXX}[1]{ % Auxiliar - no usar directamente
    \def\ai{#1}
    \ifthenelse{\equal{\etiqueta}{}}{}{
        \draw (\x, {\y + 0.5}) node[anchor=east]{\etiqueta};
    }
    \filldraw[fill=white]
        (\x, \y) rectangle +(1, 1) +(0.5, 0.5) node{\axvi}
        ++(1, 0) rectangle +(1, 1) +(0.5, 0.5) node{\axv}
        ++(1, 0) rectangle +(1, 1) +(0.5, 0.5) node{\axiv}
        ++(1, 0) rectangle +(1, 1) +(0.5, 0.5) node{\axiii}
        ++(1, 0) rectangle +(1, 1) +(0.5, 0.5) node{\axii}
        ++(1, 0) rectangle +(1, 1) +(0.5, 0.5) node{\axi}
        ++(1, 0) rectangle +(1, 1) +(0.5, 0.5) node{\ax}
        ++(1, 0) rectangle +(1, 1) +(0.5, 0.5) node{\aix}
        ++(1, 0) rectangle +(1, 1) +(0.5, 0.5) node{\aviii}
        ++(1, 0) rectangle +(1, 1) +(0.5, 0.5) node{\avii}
        ++(1, 0) rectangle +(1, 1) +(0.5, 0.5) node{\avi}
        ++(1, 0) rectangle +(1, 1) +(0.5, 0.5) node{\av}
        ++(1, 0) rectangle +(1, 1) +(0.5, 0.5) node{\aiv}
        ++(1, 0) rectangle +(1, 1) +(0.5, 0.5) node{\aiii}
        ++(1, 0) rectangle +(1, 1) +(0.5, 0.5) node{\aii}
        ++(1, 0) rectangle +(1, 1) +(0.5, 0.5) node{\ai};
}
%%%%%%%%%% Macros de tikz - Fin %%%%%%%%%%


%%%%%%%%%% Macros misceláneos - Inicio %%%%%%%%%%
\newcommand{\xmm}[1]{\texttt{XMM#1}}
\newcommand{\rax}{\texttt{RAX}}
\newcommand{\rbx}{\texttt{RBX}}
\newcommand{\rcx}{\texttt{RCX}}
\newcommand{\rdx}{\texttt{RDX}}
\newcommand{\rbp}{\texttt{RBP}}
\newcommand{\rsp}{\texttt{RSP}}
\newcommand{\reg}[1]{\texttt{R#1}}
\newcommand{\asm}[1]{\texttt{\uppercase{#1}}}
%%%%%%%%%% Macros misceláneos - Fin %%%%%%%%%%


\begin{document}


%%%%%%%%%%%%%%%%%%%%%%%%%%%%%%%%%%%%%%%%%%%%%%%%%%%%%%%%%%%%%%%%%%%%%%%%%%%%%%%
%% Carátula                                                                  %%
%%%%%%%%%%%%%%%%%%%%%%%%%%%%%%%%%%%%%%%%%%%%%%%%%%%%%%%%%%%%%%%%%%%%%%%%%%%%%%%


\thispagestyle{caratula}

\begin{center}

\includegraphics[height=2cm]{DC.png} 
\hfill
\includegraphics[height=2cm]{UBA.jpg} 

\vspace{2cm}

Departamento de Computación,\\
Facultad de Ciencias Exactas y Naturales,\\
Universidad de Buenos Aires

\vspace{4cm}

\begin{Huge}
Trabajo Práctico 2
\end{Huge}

\vspace{0.5cm}

\begin{Large}
Organización del Computador II
\end{Large}

\vspace{1cm}

Primer Cuatrimestre de 2013

\vspace{4cm}

Grupo: \textbf{Panceta y Mozzarella}

\vspace{0.5cm}

\begin{tabular}{|c|c|c|}
\hline
Apellido y Nombre & LU & E-mail\\
\hline
Laouen Louan Mayal Belloli  & 134/11 & lao.facu@gmail.com\\
Leandro Lovisolo      		& 645/11 & leandro@leandro.me\\
Lautaro José Petaccio 		& 443/11 & lausuper@gmail.com\\
\hline
\end{tabular}

\end{center}

\newpage


%%%%%%%%%%%%%%%%%%%%%%%%%%%%%%%%%%%%%%%%%%%%%%%%%%%%%%%%%%%%%%%%%%%%%%%%%%%%%%%
%% Índice                                                                    %%
%%%%%%%%%%%%%%%%%%%%%%%%%%%%%%%%%%%%%%%%%%%%%%%%%%%%%%%%%%%%%%%%%%%%%%%%%%%%%%%


\tableofcontents

\newpage


%%%%%%%%%%%%%%%%%%%%%%%%%%%%%%%%%%%%%%%%%%%%%%%%%%%%%%%%%%%%%%%%%%%%%%%%%%%%%%%
%% Introducción                                                              %%
%%%%%%%%%%%%%%%%%%%%%%%%%%%%%%%%%%%%%%%%%%%%%%%%%%%%%%%%%%%%%%%%%%%%%%%%%%%%%%%


\section{Introducción}

El objetivo de este trabajo es experimentar con el set de instrucciones SIMD de la arquitectura IA-32 de Intel. 

Con este propósito, implementamos una serie de filtros gráficos que operan sobre imágenes RGB y en escala de grises. Para cada filtro escribimos una implementación de referencia en lenguaje C, sin ningún tipo de optimización, y otra implementación en lenguaje assembler haciendo uso de las instrucciones SIMD siempre que fuera posible. Finalmente, evaluamos el rendimiento de ambas implementaciones de cada filtro, y discutimos las mejoras observadas.

En la siguiente sección analizamos cada una de las implementaciones y describimos en detalle las instrucciones SIMD usadas en cada caso.


%%%%%%%%%%%%%%%%%%%%%%%%%%%%%%%%%%%%%%%%%%%%%%%%%%%%%%%%%%%%%%%%%%%%%%%%%%%%%%%
%% Desarrollo                                                                %%
%%%%%%%%%%%%%%%%%%%%%%%%%%%%%%%%%%%%%%%%%%%%%%%%%%%%%%%%%%%%%%%%%%%%%%%%%%%%%%%


\section{Desarrollo}

\subsection{Recortar}

La salida del filtro recortar consiste en, dada una longitud $n$, recortar las esquinas de la imagen
original en forma de cuadrados de lado $n$ y unirlos para obtener la imagen destino, ubicando cada
cuadrado en su posición opuesta respecto de la imagen original.

\subsubsection{Descripción del ciclo}

Se ilustra a continuación el ciclo de la implementación C:

\begin{algorithmic}
    \FOR{$y = 1$ to $n$}
        \FOR{$x = 1$ to $n$}
            \STATE $dst[alto - n + y][ancho - n + x]$ $\leftarrow$ $src[y][x]$ \COMMENT{copio esquina superior izquierda}
            \STATE $dst[alto - n + y][x]$ $\leftarrow$ $src[y][ancho - n + x]$ \COMMENT{copio esquina superior derecha}
            \STATE $dst[y][ancho - n + x]$ $\leftarrow$ $src[alto - n + y][x]$ \COMMENT{copio esquina inferior izquierda}
            \STATE $dst[y][x]$ $\leftarrow$ $src[alto - n + y][ancho - n + x]$ \COMMENT{copio esquina inferior derecha}
        \ENDFOR
    \ENDFOR
\end{algorithmic}

La implementación assembler difiere de la anterior en cuanto a que se copia la imagen de a 16 bytes por vez,
utilizando la instrucción \texttt{MOVDQU} para cargar 16 bytes de memoria contiguos en un registro XMM 
y luego volcar el contenido del mismo en la posición de memoria correspondiente.

Además, se realiza un ajuste sobre la variable $x$ para evitar leer y escribir posiciones de memoria fuera
de la imagen y/o del cuadrado correspondiente: cuando la diferencia $|n - x|$ es menor a 16, se le resta
esta diferencia a $x$, de manera que el siguiente ciclo opere exactamente con los 16 últimos bytes de la fila actual.

El pseudocódigo a continuación describe el ciclo completo de la implementación assembler:

\begin{algorithmic}
    \FOR{$y = 1$ to $n$}
        \FOR{$x = 1$ to $n$}
            \IF{$|n - x| < 16$}
                \STATE $x$ $\leftarrow$ $x - |n - x|$
            \ENDIF
            \STATE \texttt{MOVDQU XMM0}, $src[y][x]$ \COMMENT{copio esquina superior izquierda}
            \STATE \texttt{MOVDQU} $dst[alto - n + y][ancho - n + x]$, \texttt{XMM0} \COMMENT{pego esquina superior izquierda}            
            \STATE idem demás esquinas
        \ENDFOR
    \ENDFOR
\end{algorithmic}

\subsubsection{Rendimiento}

Se observan las siguientes cantidades de ciclos y ticks de reloj al realizar 1000 iteraciones de ambas implementaciones con una imagen cuadrada de lado 512 y tamaño de esquina 100.

\begin{center}
    \begin{tabular}{|l|l|l|l|}
        \hline
        Medición & Implementación C & Implementación assembler & Relación \\
        \hline
        Ticks    & 385953030        & 23853906                 & $6.43\%$ \\
        Ciclos   & 378948.344       & 24397.717                & $6.43\%$ \\
        \hline
    \end{tabular}
\end{center}

Dado que la implementación assembler accede a la memoria de a 16 bytes por vez, la cantidad de
accesos totales será aproximadamente $\frac{1}{16} = 0.0625 = 6.25\%$ la cantidad de
accesos a memoria de la implementación C (no es exactamente esa cantidad pues los ajustes de
fin de línea introducen un pequeño adicional de lecturas, dependiendo del ancho de la imagen.)
Notemos que este número se condice con las mediciones de rendimiento realizadas.

\subsection{Halftone}

La idea de utilizar assembler fue principalmente aprobechar mejor las operaciones existentes de SIMD, teniendo en cuenta que en un registro xmm de 128 bits entran 16 bytes y que cada pixel de una imagen en escala de grises ocupa exactamente un byte; podriamos pensar en trabajar con 16 pixeles de la imagen de manera simultanea. El primer problema con el que nos encontramos, es que cada pixel puede tener un número del 0 al 255 y como en la implementación de Halftone se necesita obtener la suma de 4 pixeles, entonces podemos encontrarnos en ciertos casos con que dicha suma supere el valor 255 y no lo podamos reprecentar con un solo byte. Es por esto que tuvimos que desempaquetar los 16 bytes a word, alojando la parte baja en un registo y la parte alta en otro. De esta manera, seguimos teniendo 16 pixeles de la imagen, pero ya no de forma simultanea, sino duplicando los pasos, ya que cada instrucción tiene que ser llamada dos veces, una para cada registro.
Por otro lado, como el filtro procesa bloques de 2x2 y la cantidad de registros xmm no lo permite, en cada iteración se traen, no solos los 16 bytes siguientes de la fila actual, sino los 16 bytes correspondientes de la fila siguiente.

El segundo problema con el que nos encontramos es que la imagen es una matriz de pixeles de dimensiones desconocidas, y podriamos encontrarnos con casos donde el alto o el ancho de la misma contenga una cantidad impar de pixeles. Como el filtro funciona con bloques de 2x2, y el enunciado dice que la imagen de destino tiene que tener dimensiones pares, lo que hacemos en estos casos, es decrementar en uno el alto y/o ancho si estos si son impares.

Por otro lado, en algunos pasos del algoritmo fue necesario la utilización de mascaras y valores predefinidos, las cuales eran siempre los mismos sin importar en que iteración estemos trabajando, es por esto que las calculamos al principio antes de empezar a recorrer la imagen, para de esta manera ahorrar llamado a memoria los cuales reducen significativamente la velocidad de procesamiento de la aplicación.

estos valores son:
\begin{itemize}
  \item xmm12 $\Rightarrow$ Contiene 16 bytes packed con el mínimo
  \item xmm11 $\Rightarrow$ Contiene el mínimo packed en words
  \item xmm5 $\Rightarrow$ Contiene el máximo packed en words
  \item xmm6 $\Rightarrow$ Contiene la representación flotante de Q packed
  	\item XMM9 $\Rightarrow$ contiene en sus dos quadword el valor 0xffffffffffffffff
	\item XMM10 $\Rightarrow$ contiene en sus dos quadword el valor 0x00ff00ff00ff00ff
	\item XMM11 $\Rightarrow$ contiene en sus dos quadword el valor 0xff00ff00ff00ff00
	\item XMM12 $\Rightarrow$ 205,205,205,205,205,205,205,205
	\item XMM13 $\Rightarrow$ 410,410,410,410,410,410,410,410
	\item XMM14 $\Rightarrow$ 615,615,615,615,615,615,615,615
	\item XMM15 $\Rightarrow$ 820,820,820,820,820,820,820,820 
\end{itemize}

\subsubsection{Descripción del ciclo:}
Teniendo en cuenta que debo procesar 
El primer paso realizado dentro de un ciclo es obtener los proximos 16 pixeles de la imagen en el registro xmm

\subsection{Umbralizar}

El filtro umbralizar consiste en evaluar cada pixel de la imagen y asignarle un nuevo valor según tres criterio.
\begin{itemize}
\item Si el pixel supera el máximo pasado por parámetro a la función, se le coloca un 255.
\item Si el pixel es menor que el mínimo, también pasado por parámetro, se le coloca un 0.
\item Si el pixel está entre ($min \leq pixel \leq max$), se le asigna $\lfloor pixel/q \rfloor * q$, donde q es un parámetro.
\end{itemize}

\subsubsection{Descripción del ciclo}
El filtro se divide en cinco etapas:
\begin{enumerate}
\item \textbf{Pre-ciclo:} Se crean ciertos registros que serán de utilidad en el ciclo.
\item \textbf{Inicio del ciclo:} Puesta en 0 del registro acumulador y obtención de máscara de mínimos.
\item \textbf{Obtención de la máscara para pixeles mayores al máximo y aplicación:} Se consigue la máscara para los pixeles que superen al mayor y se aplica al acumulador.
\item \textbf{Creación de la máscara para los pixeles ($min \leq pixel \leq max$):} Se arma la máscara a partir de una nueva comparación y máscaras anteriores.
\item \textbf{Aplicación de la máscara para ($min \leq pixel \leq max$) y fin del ciclo:} Se realizan los cálculos pertinentes a los pixeles que entran en esta categoría y lo aplica a la máscara.
\end{enumerate}

\subsubsection{Pre-ciclo}
Para optimizar el procesamiento dentro del ciclo, se calcula y se guarda en registros xmm al inicio del programa:
\begin{itemize}
  \item xmm12 $\Rightarrow$ Contiene 16 bytes packed con el mínimo
  \item xmm11 $\Rightarrow$ Contiene el mínimo packed en words
  \item xmm5 $\Rightarrow$ Contiene el máximo packed en words
  \item xmm6 $\Rightarrow$ Contiene la representación flotante de Q packed
\end{itemize}

También se calcula y se guarda en el registro rcx, la cantidad de pixeles que tiene la imagen, información que se utilizará para conocer cuántos pixeles faltan por procesar en cada iteración.

\subsubsection{Inicio del ciclo}
Al comienzo del ciclo se ponen en 0 los bytes del registro xmm8 que servirá como acumulador de los nuevos valores que tendrán los pixeles.

Se leen 16 bytes contiguos desde la imagen fuente en el registro xmm1 utilizando la instrucción movdqu, se busca cuáles son iguales al mínimo (aprovechando la instrucción pcmpeqb que realiza la comparación en los 16 bytes simultaneamente) y se guarda la máscara obtenida que se usará mas tarde.

\subsubsection{Obtención de la máscara para pixeles mayores al máximo y aplicación}

Queriendo explotar al máximo las instrucciones SIMD, tratando de procesar en simultaneo 16 bytes, encontramos que esto no era posible debido a que el set de instrucciones no contempla comparaciones de greater, lower o derivadas para bytes sin signo (necesario ya que los pixeles en grayscale van del 0 al 255), por lo que fue necesario extender los bytes mediante un desempaquetado, convirtiendolos en words, para luego hacer las comparaciones correspondientes.

Se desempaquetan los 16 bytes en parte baja (primeros 8 bytes) y en parte alta convirtiendolos a words por medio de la instrucción punpcklbw.

\begin{center}
  \begin{tikzpicture}[scale=0.75]
    \registroDieciseis{\xmm{1}}{0}{6}
                      {A15}{A14}{A13}{A12}{A11}{A10}{A9}{A8}{A7}
                      {A6}{A5}{A4}{A3}{A2}{A1}{A0}

    \draw [decoration={brace,amplitude=0.25cm},decorate, thick] ( 7.75, 5.75) -- (0.25, 5.75);
    \draw [decoration={brace,amplitude=0.25cm},decorate, thick] (15.75, 5.75) -- (8.25, 5.75);

    \draw [->, thick] (12, 5.5)  --  +(0, -1.25);
    \draw [->, thick] (4, 5.5)   -- ++(0, -0.75) -- ++(-6, 0) --
                    ++(0, -2.75) -- ++(6, 0)     --  +(0, -0.75);

    \registroDieciseis{\xmm{1}}{0}{3}
                      {0}{A7}{0}{A6}{0}{A5}{0}{A4}{0}{A3}{0}{A2}{0}{A1}{0}{A0}

    \registroDieciseis{\xmm{2}}{0}{0}
                      {0}{A15}{0}{A14}{0}{A13}{0}{A12}{0}{A11}{0}{A10}{0}{A9}{0}{A8}
  \end{tikzpicture}
\end{center}

Se buscan los números que superen al máximo comparando la parte alta y baja con el registro xmm5 preparado al inicio del programa, el cuál contiene el máximo en words empaquetadas. Utilizando pcmpgtw tanto en la parte baja como la alta, obtenemos una máscara que contiene en 0xFFFF las words mayores al máximo y en 0x0000 las demás. Empaquetamos, utilizando la instrucción packsswb, la máscara relacionada a la parte baja y a la alta, dejando en 255 sólo los bytes que sean mayores al máximo para luego sumarlos en el acumulador utilizando paddusb.

\begin{center}
  \begin{tikzpicture}[scale=0.75]
    \draw (-1, 15) node[anchor=west] {Estado inicial de los registros:};

    \registroDieciseis{\xmm{5}}{0}{13}
                      {0}{Max}{0}{Max}{0}{Max}{0}{Max}{0}{Max}{0}{Max}{0}{Max}{0}{Max}

    \registroDieciseis{\xmm{4}}{0}{11}
                      {0}{A15}{0}{A14}{0}{A13}{0}{A12}{0}{A11}{0}{A10}{0}{A9}{0}{A8}

    \registroDieciseis{\xmm{3}}{0}{9}
                      {0}{A7}{0}{A6}{0}{A5}{0}{A4}{0}{A3}{0}{A2}{0}{A1}{0}{A0}

    \draw (-1, 8) node[anchor=west]
      {Luego de ejecutar \texttt{PCMPGTW} \xmm{3}, \xmm{5} y \texttt{PCMPGTW} \xmm{4}, \xmm{5}:};

    \registroDieciseis{\xmm{4}}{-1}{6}
                      {0}{0}{0}{0}{0}{0}{0}{0}{0}{0}{0}{0}{FF}{FF}{FF}{FF}

    \registroDieciseis{\xmm{3}}{1}{4}
                      {0}{0}{0}{0}{FF}{FF}{FF}{FF}{0}{0}{FF}{FF}{0}{0}{FF}{FF}

    \draw [->, thick] (12, 3.75)   --  +(0, -1.75);

    \draw [->, thick] (-0.5, 5.75) -- ++(0, -3.25) -- ++ (4.5, 0) -- + (0, -0.5);

    \draw [decoration={brace,amplitude=0.25cm},decorate, thick] (0.25, 1.25) -- ( 7.75, 1.25);
    \draw [decoration={brace,amplitude=0.25cm},decorate, thick] (8.25, 1.25) -- (15.75, 1.25);

    \draw (-1, 3) node[anchor=west, fill=white]
      {Luego de ejecutar \texttt{PACKSSWB} \xmm{3}, \xmm{4}:};

    \registroDieciseis{\xmm{3}}{0}{0}
                      {0}{0}{0}{0}{0}{0}{FF}{FF}{0}{0}{FF}{FF}{0}{FF}{0}{FF}

  \end{tikzpicture}
\end{center}



\subsubsection{Creación de la máscara para los pixeles ($min \leq pixel \leq max$)}
Al utilizar un acumulador incialmente en 0 se pudo ahorrar el paso de comparar los pixeles menores al mínimo ya que los pixeles que cumplieran esta propiedad tendrían 0 como valor.

Para conseguir la máscara, se aprovechó las máscaras previamente obtenidas (mayores al máximo e iguales al mínimo), por lo que sólo fue necesario buscar los números mayores al mínimo (haciendo un procedimiento similar al de la búsqueda del máximo), agregarle los números iguales al mínimo (mediante un por) con la máscara creada al principio del ciclo y luego sacarle los mayores al máximo (mediante un pxor) dejandome sólo los pixeles que cumplen esta condición.

\subsubsection{Aplicación de la máscara para ($min \leq pixel \leq max$) y fin del ciclo}
Al ser necesaria un división y un truncamiento para el caso, se utilizaron single precision floats.

Fue necesario desempaquetar aún más los pixeles, utilizando punpcklwd, para poder obtener los valores de estos en double words, y poder convertirlos a single precision floats mediante la instrucción cvtdq2ps. Esta es precisión suficiente para los cálculos y además brinda la posibilidad realizar más operaciones sobre los pixeles simultanteamente que con doubles.

Se realiza el desempaquetado de la parte baja anteriormente desempaquetada (xmm1) y se los convierte a single presicion floats. Luego se los divide, utilizando divps, por el registro xmm6 el cuál contiene el valor Q empaquetado en single presicion floats calculado al inicio del programa. Se lo trunca utilizando cvttps2dq, para realizar la función floor ($\lfloor \rfloor$) convirtiendose en entero, se lo convierte otra vez a float y se lo termina multiplicando, utilizando mulps, por xmm6 otra vez para finalmente ser convertido a entero, obteniendo el valor correspondiente para cáda pixel. Se vuelve a empaquetar, utilizando la instrucción packusdw, obteniendose otra vez los valores nuevos de los pixeles en words y se repite la operación con la parte alta del desempaquetado inicial (xmm2).

Al finalizar, se empaquetan los dos registros resultantes, mediante packuswb, para luego poder aplicarle la máscara anteriormente creada para el caso utilizando un pand y sumar simultaneamente con paddusb el valor obtenido al acumulador, dejando el valor correspondiente en los pixeles que cumplían con el caso.

Por último, se copian los nuevos valores de los 16 bytes en la imagen destino utilizando movdqu y se incrementan los punteros de la imágen fuente y destino para la próxima iteración. También se resta nuestro contador de pixeles por procesar, se lo compara para ver si llegó al final, caso en el que termina la ejecución, o si quedan más e igual, caso en el que vuelve a ciclar, o si restan menos de 16 bytes por agarrar, dónde se vuelve para atrás los punteros de las imágenes para que queden 16 pixeles exactamente y se pueda realizar la última iteración.

\subsubsection{Comparación con la implementación C}
El ciclo en C hace escencialmente lo mismo que la implementación en assembler, pero esta lo hace de manera simultanea utilizando SIMD. 

Estos son los detalles de las operaciones realizadas en C para 16 bytes:
\begin{itemize}
\item 16 accesos a memoria para la lectura de cada byte y 16 para la escritura
\item Se realizan a lo sumo 32 comparaciones, (2 por cáda pixel), en dónde se necesita calcular la posición del pixel cáda vez que se lo quiere leer y se realiza un acceso a memoria para esto.
\item En caso de haber caído dentro del caso ($min \leq pixel \leq max$), se realizan cáda una de las operaciones de cálculo de manera secuencial.
\end{itemize}
Detalles de las operaciones realizadas por la implementación de assembler en 16 bytes:
\begin{itemize}
\item Única lectura y escritura en memoria de 16 bytes
\item Se realizan 2 comparaciones para el caso de los pixeles mayores al máximo y para calcular los píxeles mayores al mínimo por cada 16 bytes, ya que se comparan de a 8 bytes simultaneamente utilizando SIMD. También se realiza 1 comparación cada estos 16 bytes para obtener los iguales al mínimo.
\item Si bien siempre se realiza el cálculo en punto flotante relacionado a los valores de los pixeles del caso ($min \leq pixel \leq max$), estos se realizan con 4 instrucciones SIMD, pudiendo calcular el valor de 16 bytes por ciclo.
\end{itemize}

\subsubsection{Rendimiento}
Observamos las siguientes cantidades de ciclos y ticks de reloj al realizar 100 iteraciones de ambas implementaciones con una imagen cuadrada de lado 512 y utilando los parámetros: mín = 64, máx = 128 y Q = 16.
\begin{center}
    \begin{tabular}{|l|l|l|l|}
        \hline
        Medición & Implementación C & Implementación assembler & Relación \\
        \hline
        Ticks    & 785867600      & 70172072               & $8.92\%$ \\
        Ciclos   & 7858676        & 701721                & $8.92\%$ \\
        \hline
    \end{tabular}
\end{center}

Para hacer un análisis fino, utilizamos la herramienta objdump (objdump -d -M intel -S umbralizar\_c.o) para obtener el resultado de la compilación mediante gcc.

Destacamos de este análisis el uso contínuo de variables locales (almacenadas en el stack) y el cálculo de la posición en memoria y el acceso a esta cáda vez que se quería realizar una comparación, algo evitable usando registros en ensamblador y que seguro impactaron en la performance.

\begin{verbatim}
if(src_matrix[y][x] < min) {
;Instrucciones para settear la posicion del byte
mov    eax,DWORD PTR [rbp-0x4c]
movsxd rdx,eax
movsxd rax,ebx
imul   rdx,rax
mov    rax,QWORD PTR [rbp-0x38]
add    rdx,rax
mov    eax,DWORD PTR [rbp-0x48]
cdqe   
movzx  eax,BYTE PTR [rdx+rax*1]
;Termina el setteado de la posicion del byte
cmp    al,BYTE PTR [rbp-0x70] ;Realiza la comparacion
jae    c3 <umbralizar_c+0xc3>
\end{verbatim}

Cabe destacar que el lenguaje ensamblador, además del uso de las instrucciones SIMD, que nos permiten realizar cálculos simultaneamente otorgandonos una gran ventaja contra el código en C, también nos permite explotar los recursos al máximo, pudiendo obviar cálculos de más guardandolos para su posterior uso y utilizar registros que implican menos cantidad de accesos a memoria.

Estas ventajas, sumadas a los pocos accesos de memoria al poder escribir y leer de a 16 bytes son la causa del bajo tiempo de ejecución del filtro en ASM.

\subsection{Colorizar}

\newcommand{\registro}[4]{
    \ifthenelse{\equal{#1}{}}{}{
        \draw (#2, {#3 + 0.5}) node[anchor=east]{#1};
    }
    
    \foreach \x/\byte in {#4}
    {
        \filldraw[fill=white] ({\x + #2}, #3) rectangle +(1,1) +(0.5, 0.5) node{\byte};
    }    
}

\newcommand{\registroCuatro}[7]{
    \draw (#2, {#3 + 0.5}) node[anchor=east]{#1};
    \draw   (#2, #3) rectangle +(4, 1) +(2, 0.5) node{#4}
          ++(4, 0)   rectangle +(4, 1) +(2, 0.5) node{#5}
          ++(4, 0)   rectangle +(4, 1) +(2, 0.5) node{#6}
          ++(4, 0)   rectangle +(4, 1) +(2, 0.5) node{#7};            
}


El filtro colorizar modifica cada pixel de una imagen color de acuerdo a una ecuación que depende
de los máximos valores de los canales R, G y B entre los píxeles del cuadrado de lado 3 formado
por el pixel original y sus 8 vecinos inmediatos.

El problema principal a resolver es minimizar la cantidad de accesos a memoria y paralelizar lo más
posible las comparaciones necesarias para hallar el máximo valor de cada canal en el cuadrado
de lado 3 correspondiente a cada pixel.


%%%%%%%%%%%%%%%%%%%%%%%%%%%%%%%%%%%%%%%%%%%%%%%%%%%%%%%%%%%%%%%%%%%%%%%%%%%%%%%
%% Descripción del ciclo                                                     %%
%%%%%%%%%%%%%%%%%%%%%%%%%%%%%%%%%%%%%%%%%%%%%%%%%%%%%%%%%%%%%%%%%%%%%%%%%%%%%%%


\subsubsection{Descripción del ciclo}

Para cada pixel, dividimos el ciclo en cuatro etapas:

\begin{enumerate}
    \item \textbf{Lectura de píxeles vecinos:}
        Cargamos los píxeles de cada fila del cuadrado en \xmm{1}, \xmm{2} y \xmm{3}.
    \item \textbf{Búsqueda de máximos por canal:}
        Separamos cada pixel por canal, y comparamos sucesivamente los valores de cada canal
        entre sí para hallar los máximos.
    \item \textbf{Hallar $\Phi_{R,G,B}$:}
        Hallamos los valores de $\Phi_R$, $\Phi_G$ y $\Phi_B$ utilizando los máximos
        obtenidos en la etapa anterior.
    \item \textbf{Escritura de pixel destino:}
        Computamos el valor del pixel destino y lo escribimos en la imagen.
\end{enumerate}

Veamos a continuación cada etapa en detalle.


%%%%%%%%%%%%%%%%%%%%%%%%%%%%%%%%%%%%%%%%%%%%%%%%%%%%%%%%%%%%%%%%%%%%%%%%%%%%%%%
%% Lectura de píxeles                                                        %%
%%%%%%%%%%%%%%%%%%%%%%%%%%%%%%%%%%%%%%%%%%%%%%%%%%%%%%%%%%%%%%%%%%%%%%%%%%%%%%%


\subsubsection{Lectura de píxeles vecinos}

Cargamos los píxeles de cada fila del cuadrado de lado 3 como se indica a continuación:

\begin{pseudocodigo}
    \STATE \texttt{MOVDQU} \xmm{1}, $src[y-1][x-3]$
    \STATE \texttt{MOVDQU} \xmm{2}, $src[y][x-3]$
    \STATE \texttt{MOVDQU} \xmm{3}, $src[y+1][x-3]$        
\end{pseudocodigo}

Contenido de los registros \xmm{1}, \xmm{2} y \xmm{3} en este instante:

\begin{center}
    \begin{tikzpicture}[scale=0.75]
        \registro{}{0}{0}
                 {0/$\ast$, 1/$\ast$, 2/$\ast$, 3/$\ast$, 4/$\ast$, 5/$\ast$, 6/$\ast$,
                  7/$R$,  8/$G$,  9/$B$,
                  10/$R$, 11/$G$, 12/$B$,
                  13/$R$, 14/$G$, 15/$B$}
    \end{tikzpicture}
\end{center}

El símbolo $\ast$ representa información de píxeles adyacentes que no nos interesa.

En el caso de estar leyendo los últimos píxeles de la última fila, las lecturas
se realizan desplazándonos hacia atrás sobre el eje X tantos píxeles como sea
necesario, de forma tal que la parte alta de cada registro XMM esté alineada
con el último pixel. En este caso el contenido de cada registro sería
el siguiente:

\begin{center}
    \begin{tikzpicture}[scale=0.75]
        \registro{}{0}{0}
                 {0/$R$, 1/$G$, 2/$B$,
                  3/$R$, 4/$G$, 5/$B$,
                  6/$R$, 7/$G$, 8/$B$,
                  9/$\ast$, 10/$\ast$, 11/$\ast$, 12/$\ast$, 13/$\ast$, 14/$\ast$, 15/$\ast$}
    \end{tikzpicture}
\end{center}

Luego de realizar esta lectura, reacomodamos los bytes con \texttt{PSHUFB}
para tenerlos en el mismo orden que en el caso normal, y procedemos normalmente.

%%%%%%%%%%%%%%%%%%%%%%%%%%%%%%%%%%%%%%%%%%%%%%%%%%%%%%%%%%%%%%%%%%%%%%%%%%%%%%%
%% Búsqueda de máximos                                                       %%
%%%%%%%%%%%%%%%%%%%%%%%%%%%%%%%%%%%%%%%%%%%%%%%%%%%%%%%%%%%%%%%%%%%%%%%%%%%%%%%


\subsubsection{Búsqueda de máximos por canal}

El siguiente paso es comparar los valores de cada canal entre sí para hallar los máximos.

Lo primero que hacemos es reordenar los bytes en \xmm{1}, \xmm{2} y \xmm{3}:

\begin{pseudocodigo}
    \STATE \texttt{PSHUFB} \xmm{1}, \textit{máscara}
    \STATE \texttt{PSHUFB} \xmm{2}, \textit{máscara}
    \STATE \texttt{PSHUFB} \xmm{3}, \textit{máscara}
\end{pseudocodigo}

La \textit{máscara} provista es tal que los bytes de los registros quedan
reordeados de la siguiente forma:

\begin{center}
    \begin{tikzpicture}[scale=0.75]
        \registro{}{0}{0}
                 {0/0,  1/0,    2/0,    3/0,
                  4/0,  5/$R$,  6/$R$,  7/$R$,
                  8/0,  9/$G$,  10/$G$, 11/$G$,
                  12/0, 13/$B$, 14/$B$, 15/$B$}        
    \end{tikzpicture}
\end{center}

A continuación, buscamos los máximos byte a byte entre \xmm{1} y \xmm{2} y luego entre
\xmm{1} y \xmm{3} de manera de obtener los máximos por canal ``columna a columna'':

\begin{pseudocodigo}
    \STATE \texttt{PMAXUB} \xmm{1}, \xmm{2}
    \STATE \texttt{PMAXUB} \xmm{1}, \xmm{3}    
\end{pseudocodigo}

La siguiente ilustración muestra lo que ocurre con los primeros 3 bytes de
\xmm{1}, \xmm{2} y \xmm{3}, que contienen los valores del canal B para los
tres píxeles de las filas 1, 2 y 3 respectivamente:

\begin{center}
    \begin{tikzpicture}[scale=0.75]
        \foreach \i in {1, ..., 3}
        {
            \draw (0,  {-(\i - 1) * 2 + 0.5}) node[anchor=east]{\xmm{\i}};
            \draw (0,  {-(\i - 1) * 2}) rectangle +(1, 1) +(0.5, 0.5) node{$\dots$};
            \draw (1,  {-(\i - 1) * 2}) rectangle +(5, 1) +(2.5, 0.5) node{$B_{{\i}c}$};
            \draw (6,  {-(\i - 1) * 2}) rectangle +(5, 1) +(2.5, 0.5) node{$B_{{\i}b}$};        
            \draw (11, {-(\i - 1) * 2}) rectangle +(5, 1) +(2.5, 0.5) node{$B_{{\i}a}$};
        }

        \draw (0, -5) node[anchor=west]
            {Luego de ejecutar \texttt{PMAXUB} \xmm{1}, \xmm{2} y
                               \texttt{PMAXUB} \xmm{1}, \xmm{3}:};

        \draw (0, {-7 + 0.5}) node[anchor=east]{\xmm{1}};
        \draw (0,  -7) rectangle +(1, 1) +(0.5, 0.5) node{$\dots$};
        \draw (1,  -7) rectangle +(5, 1) +(2.5, 0.5) node{$max(B_{1c}, B_{2c}, B_{3c})$};
        \draw (6,  -7) rectangle +(5, 1) +(2.5, 0.5) node{$max(B_{1b}, B_{2b}, B_{3b})$};        
        \draw (11, -7) rectangle +(5, 1) +(2.5, 0.5) node{$max(B_{1a}, B_{2a}, B_{3a})$};        
    \end{tikzpicture}
\end{center}

Lo mismo ocurre con el resto de los bytes que contienen los valores para los canales R y G.

En este instante el contenido de \xmm{1} es el siguiente:

\begin{center}
    \begin{tikzpicture}[scale=0.75]
        \registro{\xmm{1}}{0}{0}
                 {0/0,  1/0,      2/0,      3/0,
                  4/0,  5/$R_c$,  6/$R_b$,  7/$R_a$,
                  8/0,  9/$G_c$,  10/$G_b$, 11/$G_a$,
                  12/0, 13/$B_c$, 14/$B_b$, 15/$B_a$}
    \end{tikzpicture}
\end{center}

Donde $B_a = max(B_{1a}, B_{2a}, B_{3a})$, $\dots$, $R_c = max(R_{1c}, R_{2c}, R_{3c})$.

Para encontrar $max_R$, $max_G$ y $max_B$ aún nos falta maximizar las tuplas
$(R_a, R_b, R_c)$, $(G_a, G_b, G_c)$ y $(B_a, B_b, B_c)$.

Extraemos las tres triplas $(R, G, B)$ del registro \xmm{1} usando
\texttt{PSHUFB} con las máscaras correspondientes y las almacenamos por separado en
\xmm{1}, \xmm{2} y \xmm{3}, de manera de tener un único valor R, G y B en cada uno de
estos registros, para poder buscar los máximos en paralelo columna a columna con
\texttt{PMAXUB}, como hicimos antes.

Gráficamente, podemos representar esta serie de pasos de la siguiente manera:

\begin{center}
    \begin{tikzpicture}[scale=0.75, >=stealth]
        \registro{\xmm{1}}{0}{10}
                 {0/0,  1/0,      2/0,      3/0,
                  4/0,  5/$R_c$,  6/$R_b$,  7/$R_a$,
                  8/0,  9/$G_c$,  10/$G_b$, 11/$G_a$,
                  12/0, 13/$B_c$, 14/$B_b$, 15/$B_a$}

        \draw [->, thick] (15.5, 10) -- (15.5, 8); % Ba
        \draw [->, thick] (14.5, 10) -- (15.5, 6); % Bb
        \draw [->, thick] (13.5, 10) -- (15.5, 4); % Bc

        \draw [->, thick] (11.5, 10) -- (11.5, 8); % Ga
        \draw [->, thick] (10.5, 10) -- (11.5, 6); % Gb
        \draw [->, thick]  (9.5, 10) -- (11.5, 4); % Gc

        \draw [->, thick]  (7.5, 10) --  (7.5, 8); % Ra
        \draw [->, thick]  (6.5, 10) --  (7.5, 6); % Rb
        \draw [->, thick]  (5.5, 10) --  (7.5, 4); % Rc        

        \draw (0, 9) node[anchor=west, fill=white]{Luego de extraer las triplas:};

        \registro{\xmm{1}}{0}{7}
                 {0/0,  1/0,  2/0,  3/0,
                  4/0,  5/0,  6/0,  7/$R_a$,
                  8/0,  9/0,  10/0, 11/$G_a$,
                  12/0, 13/0, 14/0, 15/$B_a$}

        \registro{\xmm{2}}{0}{5}
                 {0/0,  1/0,  2/0,  3/0,
                  4/0,  5/0,  6/0,  7/$R_b$,
                  8/0,  9/0,  10/0, 11/$G_b$,
                  12/0, 13/0, 14/0, 15/$B_b$}

        \registro{\xmm{3}}{0}{3}
                 {0/0,  1/0,  2/0,  3/0,
                  4/0,  5/0,  6/0,  7/$R_c$,
                  8/0,  9/0,  10/0, 11/$G_c$,
                  12/0, 13/0, 14/0, 15/$B_c$}

        \draw (0, 2) node[anchor=west]
            {Luego de ejecutar \texttt{PMAXUB} \xmm{1}, \xmm{2} y
                               \texttt{PMAXUB} \xmm{1}, \xmm{3}:};

        \registro{\xmm{1}}{0}{0}
                 {0/0,  1/0,  2/0,  3/0,
                  4/0,  5/0,  6/0,  7/$R^\ast$,
                  8/0,  9/0,  10/0, 11/$G^\ast$,
                  12/0, 13/0, 14/0, 15/$B^\ast$}
    \end{tikzpicture}
\end{center}

Donde $R^\ast = max_R$, $G^\ast = max_G$ y $B^\ast = max_B$.

Finalmente extraemos los máximos con \texttt{PSHUFD} y \texttt{MOVD},
y los guardamos en registros de propósito general, para usarlos en la
siguiente etapa.


%%%%%%%%%%%%%%%%%%%%%%%%%%%%%%%%%%%%%%%%%%%%%%%%%%%%%%%%%%%%%%%%%%%%%%%%%%%%%%%
%% Hallar Phi                                                               %%
%%%%%%%%%%%%%%%%%%%%%%%%%%%%%%%%%%%%%%%%%%%%%%%%%%%%%%%%%%%%%%%%%%%%%%%%%%%%%%%


\subsubsection{Hallar $\Phi_{R,G,B}$}

Procedemos a computar $\Phi_R$, $\Phi_G$ y $\Phi_B$ de forma secuencial,
luego empaquetamos en \xmm{1} los valores hallados como se indica a continuación:

\begin{center}
    \begin{tikzpicture}[scale=0.75]
        \registroCuatro{\xmm{1}}{0}{0}{$\ast$}{$\Phi_R$}{$\Phi_G$}{$\Phi_B$}
    \end{tikzpicture}
\end{center}

Empaquetamos $\Phi_{R,G,B}$ de esta manera para facilitar las
comparaciones que realizaremos en la siguiente y última etapa.


%%%%%%%%%%%%%%%%%%%%%%%%%%%%%%%%%%%%%%%%%%%%%%%%%%%%%%%%%%%%%%%%%%%%%%%%%%%%%%%
%% Pixel destino                                                             %%
%%%%%%%%%%%%%%%%%%%%%%%%%%%%%%%%%%%%%%%%%%%%%%%%%%%%%%%%%%%%%%%%%%%%%%%%%%%%%%%


\subsubsection{Escritura de pixel destino}

Antes de escribir el pixel destino, debemos hallar el valor de cada canal.
Recordemos las ecuaciones:

\begin{center}
    $I_{dst_R}(i, j) = min(255, \Phi_R \ast I_{src_R}(i, j))$ \\
    $I_{dst_G}(i, j) = min(255, \Phi_G \ast I_{src_G}(i, j))$ \\
    $I_{dst_B}(i, j) = min(255, \Phi_B \ast I_{src_B}(i, j))$
\end{center}

Procedemos a hallar $I_{src_{R,G,B}}(i, j)$:

\begin{pseudocodigo}
    \STATE \texttt{MOVDQU} \xmm{2},   $src[j][i]$      \COMMENT{leo el pixel original}
    \STATE \texttt{PSHUFB} \xmm{2},   \textit{máscara} \COMMENT{reacomodo y limpio registro}
    \STATE \texttt{CVTDQ2PS} \xmm{2}, xmm{2}           \COMMENT{convierto cada canal a float}
\end{pseudocodigo}

En este instante \xmm{2} contiene el valor de cada canal del pixel original
convertido a punto flotante como se ilustra a continuación:

\begin{center}
    \begin{tikzpicture}[scale=0.75]
        \registroCuatro{\xmm{2}}{0}{0}{$\ast$}{$I_{src_R}$}{$I_{src_G}$}{$I_{src_B}$}
    \end{tikzpicture}
\end{center}

Multiplicamos cada canal por el valor $\Phi$ correspondiente y buscamos
el mínimo entre este producto y 255:

\begin{pseudocodigo}
    \STATE \texttt{MULPS}    \xmm{2}, \xmm{1} \COMMENT{\xmm{2} $\leftarrow$ $\Phi \ast I_{src}$}
    \STATE \texttt{MOVDQU}   \xmm{1}, \textit{tupla 255}
        \COMMENT{cargo el valor 255 en los 3 floats más bajos de \xmm{1}}
    \STATE \texttt{MINPS}    \xmm{1}, \xmm{2}
        \COMMENT{\xmm{1} $\leftarrow$ $min(255, \Phi \ast I_{src}) = I_{dst}$}
    \STATE \texttt{CVTPS2DQ} \xmm{1}, \xmm{1} \COMMENT{convertimos a enteros los $I_{dst}$ hallados}
\end{pseudocodigo}

El contenido de de \xmm{1} es ahora:

\begin{center}
    \begin{tikzpicture}[scale=0.75]
        \registroCuatro{\xmm{1}}{0}{0}{$\ast$}{$I_{dst_R}$}{$I_{dst_G}$}{$I_{dst_B}$}
    \end{tikzpicture}
\end{center}

Si bien los valores $I_{dst_{R,G,B}}$ resultantes de la conversión de punto flotante a
entero sin signo son de 32 bits, recordemos que nunca se exceden de 255, por lo que podemos
asumir que nuestro dato está en el byte más bajo de cada entero, mientras que en los bytes
restantes hay ceros. Entonces el contenido de \xmm{1} es:

\begin{center}
    \begin{tikzpicture}[scale=0.75]
        \registro{\xmm{1}}{0}{0}
                 {0/$\ast$,  1/$\ast$, 2/$\ast$, 3/$\ast$,
                  4/0,       5/0,      6/0,      7/$I_{dst_R}$,
                  8/0,       9/0,      10/0,     11/$I_{dst_G}$,
                  12/0,      13/0,     14/0,     15/$I_{dst_B}$}
    \end{tikzpicture}
\end{center}

Reordenamos los bytes con \texttt{PSHUFB}:

\begin{center}
    \begin{tikzpicture}[scale=0.75]
        \registro{\xmm{1}}{0}{0}
                 {0/$\ast$,  1/$\ast$, 2/$\ast$, 3/$\ast$, 4/0,
                  5/0, 6/0,  7/0, 8/0,  9/0,  10/0, 11/0, 12/0,
                  13/$I_{dst_R}$, 14/$I_{dst_G}$, 15/$I_{dst_B}$}
    \end{tikzpicture}
\end{center}

Para terminar, extraemos $I_{dst_R}$, $I_{dst_G}$ y $I_{dst_B}$ con \texttt{MOVD}
y escribimos el pixel obtenido en la imagen destino.


%%%%%%%%%%%%%%%%%%%%%%%%%%%%%%%%%%%%%%%%%%%%%%%%%%%%%%%%%%%%%%%%%%%%%%%%%%%%%%%
%% Descripción del ciclo                                                     %%
%%%%%%%%%%%%%%%%%%%%%%%%%%%%%%%%%%%%%%%%%%%%%%%%%%%%%%%%%%%%%%%%%%%%%%%%%%%%%%%


\subsubsection{Comparación con la implementación C}

Escencialmente, la versión C de este filtro realiza la misma secuencia de operaciones,
pero accediendo a memoria de a un byte por vez, y realizando las comparaciones
sin paralelismo. Esto significa que se realizan:

\begin{itemize}
    \item $9 * 3 = 27$ accesos a memoria para obtener cada canal de los píxeles vecinos,
    \item 3 accesos a memoria para leer el pixel original,
    \item 3 accesos a memoria para escribir el pixel destino, y
    \item 24 comparaciones para hallar $max_{R,G,B}$.
\end{itemize}

Esto es un total de 33 accesos a memoria y 24 comparaciones.

En contraste, la implementación assembler realiza:

\begin{itemize}
    \item 3 accesos a memoria para obtener los píxeles vecinos,
    \item 1 acceso a memoria para leer el pixel original,
    \item 1 acceso a memoria para escribir el pixel destino, y
    \item 6 comparaciones para hallar $max_{R,G,B}$.
\end{itemize}

Es decir, un total de 5 accesos a memoria y 6 comparaciones.

Expresado de forma relativa, la implementación assembler realiza
el $15.2\%$ de accesos a memoria y el $25\%$ de comparaciones que su
contraparte en lenguaje C.


%%%%%%%%%%%%%%%%%%%%%%%%%%%%%%%%%%%%%%%%%%%%%%%%%%%%%%%%%%%%%%%%%%%%%%%%%%%%%%%
%% Rendimiento                                                               %%
%%%%%%%%%%%%%%%%%%%%%%%%%%%%%%%%%%%%%%%%%%%%%%%%%%%%%%%%%%%%%%%%%%%%%%%%%%%%%%%


\subsubsection{Rendimiento}

Observamos las siguientes cantidades de ciclos y ticks de reloj al realizar 100 iteraciones de ambas implementaciones con una imagen cuadrada de lado 512 y parámetro $\alpha = 0.5$:

\begin{center}
    \begin{tabular}{|l|l|l|l|}
        \hline
        Medición & Implementación C & Implementación assembler & Relación \\
        \hline
        Ticks    & 53062000890      & 1213052436               & $2.28\%$ \\
        Ciclos   & 530620000        & 12130524                 & $2.28\%$ \\
        \hline
    \end{tabular}
\end{center}

En la sección anterior evaluamos las relaciones de accesos a memoria y comparaciones
entre ambas implementaciones. Sin realizar un mayor análisis, éstas nos sirven
para tener una idea intuitiva de las mejoras que podríamos esperar de haber optimizado
utilizando instrucciones SIMD. Sin embargo el rendimiento obtenido fue casi un orden
de magnitud superior.

Luego de analizar el código assembler generado por el compilador, identificamos una
gran cantidad de accesos a memoria adicionales, producidos por el uso de variables locales
almacenadas en la pila. Esto ocurre tanto en el ciclo principal como en varios auxiliares.
Atribuímos la diferencia de performance a estos accesos a memoria que no tuvimos en cuenta
en nuestro análisis inicial.

\subsection{Efecto plasma}

Este filtro consiste en obtener los pixeles de la imagen original y aplicarle una serie de cuentas para luego insertar el resultado en la imagen de destino. Las cuentas que se realizan son las siguientes. \\

\begin{itemize}
	\item $prof = ( x\_scale*sin\_taylor(i/8.0) + y\_scale*sin\_taylor(j/8.0) )/2$. Donde $sin\_taylor$ es una función que calcula el seno de $x$ a travez de una aproximación, y $x\_scale$,$y\_scale$ son parámetros de entrada de la función.
	\item Si el valor optenido es mayor a $255$ el resultado es $255$.
	\item Si el contenido es menor a $0$ el resultado es $0$.
	\item En otro caso el se trunca el resultado.
\end{itemize}

\subsubsection{Pseudo codigo del algoritmo en C:}

\begin{pseudocodigo}
   \FOR {alto de la imagen}
		\FOR {ancho de la imagen}
			\STATE
			\STATE $prof = ( x\_scale*sin\_taylor(i/8.0) + y\_scale*sin\_taylor(j/8.0) )/2$ 
			\STATE $newValue = prof*g\_scale + src\_matrix[i][j]$ 
			\STATE
			\IF {$newValue > 255$}
				\STATE $newValue = 255$
			\ELSE 
				\IF{$newValue < 0$}
					\STATE $newValue = 0$
				\ENDIF 
			\ENDIF
			\STATE $dst\_matrix[i][j] = floor(newValue)$ 
		\ENDFOR
	\ENDFOR
\end{pseudocodigo}

$sin\_taylor(x)$:


\begin{pseudocodigo}
   \FOR {alto de la imagen}
		\FOR {ancho de la imagen}
			\STATE
			\STATE $pi	= 3.14159265359$
			\STATE $k 			= floor(x/(2*pi))$
			\STATE $r 		= x - k*2*pi$
			\STATE $x 				= r - pi$
			\STATE $y 		= x - (pow(x,3)/6) + (pow(x,5)/120) - (pow(x,7)/5040)$
			\STATE
			\STATE devolver $y$
			\STATE
		\ENDFOR
	\ENDFOR
\end{pseudocodigo}

\subsubsection{Idea general:}
La motivación de utilizar assembler para la resolución de esta aplicación es fundamentalmente la aplicación de las herramientas SIMD para mejorar la performanse del algoritmo. Teniendo en cuenta que el mismo trabaja haciendo iteraciónes sobre los píxeles de una imagen y procesando los mismos, si podemos procesar más de un pixel por iteración mejoramos de forma significativa la performance de la aplicación.\\

La idea principal fue implementar el algoritmo utilizando los registros $xmm$ para poder procesar 16 píxeles por iteración. Ante esta idea, el problema que nos surgio fue el siguiente:\\

Como se puede ver en el pseudocodigo debemos llamar a la función $sin\_taylor(j/80.0)$ y $sin\_taylor(i/80.0)$, donde $(i,j)$ hacen referencia a la posición del pixel dentro de la matriz de la imagen. No podemos asumir nada en cuanto al valor de $j$ e $i$, incluso, estos pueden ser tan grandes como quieran, dependiendo el tamaño de la imagen, haciendo en ciertos casos que el algoritmo en assembler tenga fallas si se superara la precisión de un byte. 

El algoritmo en C esta implementado con float, lo cual es una precisión razonable para las dimensiones de imagenes con que se suele trabajar, es por esto que decidimos trabajar los datos con doubleword's, codificandolos como punto flotante de simple precisión.

Al trabajar con los datos en esta precisión lo que tenemos es que para trabajar con $16$ pixeles necesitamos $4$ registros $xmm$ para cargar los indices de las columnas a los que estos pertenecen, más otro registro para cargar el valor empaquetado del indice de la fila a la que pertencen. Si miramos la función $sin\_taylor$ veremos que necesitamos entonces, 5 registros para el acumulador de la cuenta, 5 registros para mantener el valor original mientras que en otros 5 registros se van calculando cada termino del polinomio. Esto nos da una suma de $15$ registros $xmm$. hay maneras de poder trabajar con 16 pixeles reutilizando registros y llamando a memoria, pero como justamente no queremos hacer llamados a memoria, nos parecio más eficiente procesar 8 pixeles por iteración y utilizar los registros $xmm$ sobrantes para precargar valores que seran utilizados en todas las iteraciones. \\

\subsubsection{Pre-ciclo principal:}
una vez visto esto, pasemos a nombrar todos los registros que pre-seteamos antes de empezar el ciclo, y que luego serán utilizado dentro del mismo.

estos valores son:
\begin{itemize}
	\item $xmm0$ $\Rightarrow$ $x\_scale$;$x\_scale$;$x\_scale$;$x\_scale$
	\item $xmm1$ $\Rightarrow$ $y\_scale$;$y\_scale$;$y\_scale$;$y\_scale$
	\item $xmm2$ $\Rightarrow$ $g\_scale$;$g\_scale$;$g\_scale$;$g\_scale$
	\item $xmm3$ $\Rightarrow$ $pi$;$pi$;$pi$;$pi$
	\item $xmm4$ $\Rightarrow$ $2.0$;$2.0$;$2.0$;$2.0$
	\item $xmm5$ $\Rightarrow$ $3$;$2$;$1$;$0$
	\item $xmm6$ $\Rightarrow$ $7$;$6$;$5$;$4$
	\item $xmm7$ $\Rightarrow$ $6.0$,$6.0$;$6.0$;$6.0$
	\item $xmm8$ $\Rightarrow$ $120.0$;$120.0$;$120.0$;$120.0$
	\item $xmm9$ $\Rightarrow$ $5040.0$;$5040.0$;$5040.0$;$5040.0$
\end{itemize}

\subsubsection{Descripción del ciclo:}
\begin{itemize}
	\item El primer paso realizado dentro de la iteración es guardar en $2$ registros $xmm$ los valores de los indices de los pixeles que se van a procesar, Sea $(i,j)$ el primero de los $8$ pixeles que seran procesados, los valores que se necesitan guardar son: $j$,$j+1$,$j+2$,$j+3$,$j+4$,$j+5$,$j+6$,$j+7$,$j+8$, y también el valor de $i$, el cual es el mismo para todos los pixeles ya que estamos trabajando sobre una misma fila.
	
	Para hacer esto, necesitamos los siguientes pasos:
	\begin{itemize}
	 	\item Copiamos el valor de los registros $xmm5$ y $xmm6$ a los registros $xmm10$ y $xmm11$ utilizando la instrucción MOVDQU, obteniendo en ellos los valores del $0$ al $7$.
	 	\item En un registro $xmm$ se le pone en todos sus doubleword's el valor de $j$ utilizando la instrucciones MOVQ que mueve un doubleword en la parte baja del registro $xmm$ y extiende el resto del registro con $0$'s y SHUFPS que lo expande en las demas doubleword's del registro. 
	 	\item Se lo codifica a punto flotante con la instrucción CVTDQ2PS. 
	 	\item Luego se procede a sumarle a los registros $xmm10$ y $xmm11$ el registro donde empaquetamos el valor de $j$, para esto Utilizamos la instrucción ADDPS que realiza la suma empaquetada de doubleword codificados como puntos flotantes.
	\end{itemize}
	
	(El algoritmo procesa mas adelante el valor de $i$, por lo que en esta instancia no hacemos nada con el).

	\item El segundo paso a realizar es dividir los valores de los $j$'s por $8$, pero como tenemos ya en un registro $xmm$ empaquetado el número $2.0$ como doubleword's, lo que hacemos es dividir tres veces por este registro utilizando la intrucion DIVPS.

	\subsubsection{función $Sin\_taylor$:}
	\item  El próximo paso es implementarle a estos valores la funcion $sin\_taylor$. Para esto copiamos a los registros $xmm12$ y $xmm13$ el contenido de los registros $xmm10$,$xmm11$ respectivamente, y seguimos los siguientes pasos:

	\begin{itemize}
		\item Dividimos los valores de cada pixel por $2*pi$. Como tenemos estos valores en los registros $xmm3$ y $xmm4$ utilizamos nuevamente las intrucciones DIVPS para dividir $xmm12$ y $xmm3$ por los registros anteriores, y luego aplicamos las instrucciones CVTTPS2DQ y CVTTPS2DQ para obtener la parte entera.

		\item El proximo paso de la funcion de taylor es restarle a los valores originales el resultado anterior previamente multiplicado por $2*pi$. Primero, multiplicamos utilizando la instrucción MULPS, los registros $xmm3$ y $xmm4$ a los registros $xmm12$,$xmm13$ donde estan los resultados anteriores, y luego a los registros $xmm10$ y $xmm11$ donde se encuentran los valores originales le restamos los registros $xmm12$,$xmm13$ correspondientementes utilizando la instrucción  SUBPS.\\
		En este paso hemos perdido los valores originales que teniamos empaquetados en $xmm10$ y $xmm11$ pero esto no nos importa ya que no los necesitamos más.

		\item Luego vamos a restarle a lo recien obtenido el valor de $pi$ empaquetado en el registro $xmm3$, y a continuación una seguidilla de pasos repetidos de forma casi igual, por lo que consideramos mejor una explicación general para estos pasos, y no una lectura tediosa y repetitiva de la realización de los mismos.

		Estos pasos son utilizados para obtener el valor siguiente, $Y = X - X^3/6 + X^5/120 - X^7/5040$ en cada doubleword de dos registros $xmm$ de forma empaquetada ($X$ hace referencia a los valores obtenidos mediante los pasos anteriores que estan empaquetados en los registros $xmm12$ y $xmm13$).
		Para realizar esta cuenta, se guarda el valor de $xmm12$,$xmm13$ en dos registros $xmm$ para poder salvarlos, luego se copian de vuelta estos valores en otros dos registros $xmm$ que seran utilizados como acumulador. y lo siguiente a realizar es ir multiplicando y dividiendo estos registros para obtener cada uno de los terminos del polinomio y sumarselo o restarselo a los registros acumuladores para ir teniendo el valor del polinomio en los acumuladores. Los dividendos utilizados estan ya pre-seteados en los registros $xmm7$,$xmm8$,$xmm9$ de forma empaquetada codificados en punto flotante.
	\end{itemize}

	\item Una vez que tenemos el $sin\_taylor$ para los indices $j$'s de forma empaquetada en los registros $xmm12$,$xmm13$, pasamos a empaquetar el indice $i$ en el registro $xmm10$ y realizamos los mismos pasos anteriores para obtener el $sin\_taylor$ de $i/8.0$. Con lo cual al final de los pasos conseguimos en el registro $xmm11$ el valor empaquetado en doubleword's de $sin\_taylor$ para el indice $i/8.0$, (el cual es el mismo valor para todos los pixeles por pertenecer a la misma fila).

	\item Lo siguiente, es obtener el valor de $prof(i,j)$. Para esto, primero le multiplicamos a los registro $xmm12$,$xmm13$ donde estan empaquetados los taylor's de los indices $j$'s$/8.0$ el registro $xmm1$, que es donde esta empaquetado el valor de $y\_scale$, y le multiplicamos  a $xmm11$ el registro $xmm0$ ya que en este esta empaquetado el valor de $x\_scale$. Una vez realisado esto, le sumamos a los registros $xmm12$,$xmm13$ el registro $xmm11$ y los dividimos por el registro $xmm4$ que es el que tiene el valor $2.0$ empaquetado en doubleword.

	De esta menera queda guardado en los registros $xmm12$,$xmm13$ la funcion $Prof(i,j)$ de los indices de los pixeles que se estan procesando. En $xmm12$ estan los valores para los pixeles de las pisiciones $(i,j)$ hasta $(i,j+3)$ y en $xmm13$ estan los valores para los pixeles de las posiciones $(i,j+4)$ hasta $(i,j+7)$

	\item El siguiente paso consiste en traer los valores de los pixeles de la memoria y desempaquetarlos guardando el valor en los registros $xmm14$ y $xmm15$ (en $xmm14$ se guardan los primeros $4$ pixeles y en $xmm15$ el resto).

	Teniendo estos valores hacemos la suma de $xmm12$ con $xmm14$ y $xmm13$ con $xmm15$ y tenemos el valor ya listo para ser empaquetado y guardado en el destino.

	Teniendo en cuenta que este valor puede ser mayor a $255$ o menor a $0$ para empaquetar utilizamos la instrucción PACKUSDW que empaqueta saturando integer's sin signo, dejando asi en $255$ todos aquellos valores mayores al mismo y en $0$ a los valores negativos.

\end{itemize}

\subsubsection{Comparación con el lenguaje C}

Mirando el pseudocodigo del algoritmo en C vemos para empezar que la cantidad de iteraciones es 8 veces mayor en C que en assembler al utilizar SIMD. ya que mientras que en C se procesa de a un solo pixel por vez en assembler se procede a procesar 8.

Otro aspecto a tener en cuenta es que en la implementación en assembler tenemos que obtener a partir de la posición del primer pixel los valores de los $j$'s para los otros $7$ pixeles que se van a procesar, si vien esto en C no hay que hacerlo, el costo de este paso del algoritmo esta esta compensado ya que al realizar 8 veces menos cantidad de iteraciones hay que avanzar ocho veces menos los contadores. Podemos concluir entonces que en ciertos casos, si vien paralelizar supone realizar algunos pasos extras, hay que ver si los mismos son compensados o no. pueden existir casos donde la cantidad de pasos extras a compensar sean demasiados y termine bajando la performance comparandolo con algoritmos comunes.

Tambien podemos ver que en la version C cada vez que se llama a la funcion $sin\_taylor$, se define el valor de $pi$ mientras que en la implementacion assembler el mismo ya esta predefinido antes de comenzar el ciclo, lo mismo pasa con los valores $6$, $120$ y $5040$. esto mejora la performance, pero no se lo puede considerar como una mejora de SIMD ya que lo mismo podria ser facilmente aplicado en la version C del algoritmo.\\

\underline{\textbf{Consideraciones a la hora de comparar la versión C con assembler:}}

Tengamos en cuenta para la comparación de ambas implementaciones que en C por cada pixel se hace la llamada a la funcion $sin\_taylor$ que como sabemos supone ciertos pasos para satisfacer la convencion C y pasarle los parametros a la misma, Esto lo convierten mas lento que si la misma funcion estuviera definida dentro del mismo algoritmo.

También pasa que tanto en la implementación C como en la implementación assembler, podriamos mejorarlas calculando la funcion $sin\_taylor$($i/8.0$) una sola vez por cada fila, de esta manera se reduce significativamente la llamada a la función, mejorando la performanze. En este punto es importante ver que esta mejora es $8$ veces menos significatica para assembler que para C, esto se debe a que en el mismo ya se procesan 8 pixeles por vez. Por lo que vemos que mismo en algoritmos ineficientes puede pasar que al paralelizar estemos obteniendo de por si un algoritmo mejorado y aplicable a la versión no paralelizada del mismo sin siquiera saberlo.

\subsubsection{Rendimiento}

Observamos las siguientes cantidades de ciclos y ticks de reloj al realizar 100 iteraciones de ambas implementaciones con una imagen cuadrada de lado 512.
\begin{center}
    \begin{tabular}{|l|l|l|l|}
        \hline
        Medición & Implementación C & Implementación assembler & Relación \\
        \hline
        Ticks    & 40908368256      & 960464334            & $2.35\%$ \\
        Ciclos   & 409083712        & 9604643              & $2.35\%$ \\
        \hline
    \end{tabular}
\end{center}

\subsection{Rotar}

El filtro rotar consiste en rotar los píxeles de la imagen 45 grados en sentido antihorario.

Rotar es el filtro que menos se presta para la paralelización. Esta dificultad proviene de la imposibilidad de leer bytes y procesarlos simultáneamente debido a que el valor de cada pixel en la imagen destino se obtiene calculando, según su posición, qué pixel corresponde de la imagen fuente, lo que implica traer uno a uno los píxeles de la imagen fuente para colocarlos en en lugar indicado en la imágen destino.

Analizando el algoritmo, llegamos a la conclusión de que la mejor manera de explotar la simultaneidad que nos brindan las instrucciones SIMD, es utilizarlas para calcular las posiciones de dónde sacar el pixel para colocarlo en la imagen destino. 

Inicialmente este proceso fue realizado utilizando single precision floats para poder aprovechar SIMD al máximo y realizar cálculos para conseguir las posiciones de donde sacar los píxeles a la vez. Este intento se vió frustrado al correr el filtro y compararlo contra las imágenes provistas por la cátedra. Si bien la imagen coincidía visualmente, se encontraron errores en 27 píxeles. Un análisis de estos píxeles pusieron en evidencia un problema de precisión, el cual se intentó arreglar utilizando la instrucción roundps sin conseguir mejores resultados. Se solucionó el problema realizando los cálculos usando double precision floats con SIMD para luego convertirlos a floats y finalmente convertirlos a enteros para poder funcionar como posición.


\subsubsection{Descripción del ciclo}

El filtro se divide en cinco etapas:

\begin{enumerate}
\item \textbf{Pre-ciclo principal e inicio:} Se crean ciertos registros que serán de utilidad en el ciclo.
\item \textbf{Ciclo para obtención de posiciones:} Se obtiene mediante un ciclo los valores de x e y para las siguientes cuatro posiciones.
\item \textbf{Cálculo de $u$ y $v$:} Cálculo de los valores $u$ y $v$ correspondientes a $x$ e $y$ en la imagen fuente.
\item \textbf{Ciclo de almacenamiento de los valores correspondientes a las posiciones procesadas:} Se buscan los valos de los píxeles en la imagen fuente que se escribirán en la imagen destino.
\item \textbf{Escritura de los píxeles en la imagen destino y fin del ciclo principal:}
\end{enumerate}


\subsubsection{Pre-ciclo principal}

Para optimizar la implementación, se realiza el cálculo de $cx$, $cy$ y $\sqrt[2]{2}/2$ al inicio del programa. Estos valores se almacenan en registros como doubles empaquetados utilizando \texttt{SHUFPS} para tener cada valor en double words y luego convertirlos a doubles con \texttt{CVTDQ2PD}.

También se ponen en 0 los registros \reg{11} y \reg{15} que servirán para establecer en que posición $x$ e $y$ respectivamente se encuentra en la imagen.

El ciclo principal comienza poniendo en 0 mediante la instrucción \texttt{POR} el registro \xmm{7} el cuál se encargará de acumular los píxeles hasta llenarse sus 16 bytes.


\subsubsection{Ciclo para obtención de posiciones}

Este ciclo interno se encarga de poner en dos registros temporales los valores de las posiciones $x$ e $y$ de cuatro píxeles sucesivos. Una vez realizado esto, para poder trabajar con los cuatro píxeles utilizando doubles, se mueve el contenido de estos registros a registros XMM dónde se los convierte, utilizando \texttt{CVTDQ2PD} y se los almacena en cuatro registros XMM, dos con los dos doubles más bajos y altos de las coordenadas $x$ y dos con los dos doubles más bajos y altos de las coordenadas $y$.

\begin{figure}[h!]
    \centering
    \begin{tikzpicture}[scale=0.75]
        \draw (0, 5) node[anchor=west]{Posiciones $x$:};
        
        \registroCuatro{}{0}{3}{0}{1}{2}{3}

        \draw (0, 2) node[anchor=west]{Posiciones $y$:};

        \registroCuatro{}{0}{0}{0}{0}{0}{0}
    \end{tikzpicture}
    \caption{Estado de los registros XMM asignados para $x$ e $y$ en la primera iteración del ciclo principal}
\end{figure}


\subsubsection{Cálculo de u y v}

Se opera sobre los doubles utilizando instrucciones SIMD, obteniendo los valores $u$ y $v$, es decir, los valores $x$ e $y$ respectivamente de dónde conseguir el pixel a poner en la imagen fuente.

Haciendo un análisis de las ecuaciones, decidimos realizar los cálculos de $u$ y $v$ de la siguiente manera:
\begin{itemize}
\item Dividir y calcular $u$ como $u1 = cx + (\sqrt[2]{2}/2) * (x-cx)$ y $u2 = -(\sqrt[2]{2}/2) * (y-cy)$.
\item Dividir y calcular $v$ como $v1 = cy + (\sqrt[2]{2}/2) * (x-cx)$ y $v2 = +(\sqrt[2]{2}/2) * (y-cy)$.
\end{itemize}

De esta manera se puede calcular con instrucciones SIMD, utilizando \texttt{SUBPD}, $x - cx$, multiplicándolo después con \texttt{MULPD} por $\sqrt[2]{2}/2$ y utilizando luego este resultado para formar $v1$ y $u1$ con la suma de $cy$ y $cx$ respectivamente. Por último calcular $u2$ y $v2$, $(\sqrt[2]{2}/2)*(y - cy)$, de manera similar, restándolo a $u1$ y sumándolo a $v1$.

Al finalizar el cálculo se convierten de doubles a floats utilizando \texttt{CVTPD2PS}, de floats a enteros con \texttt{CVTTPS2DQ} y se shiftean hacia la izquierda los registros con las partes altas de $x$ y de $y$ (ahora $u$ y $v$), mediante \texttt{PSLLQD}, para luego agregarles en la parte baja los valores de las partes bajas calculados utilizando la instrucción \texttt{POR}, dejando 4 double words packed con valores de $u$ y 4 double words packed con valores de $v$.

\subsubsection{Ciclo de almacenamiento de los valores correspondientes a las posiciones procesadas}

Antes de empezar este ciclo, se pone el registro \rbp en 0 mediante un \texttt{XOR} ya que servirá como acumulador temporal del ciclo.

Se pasa a un pequeño ciclo el cuál mediante \texttt{MOVD} y \texttt{PSLLDQ} (que realiza un shift a la izquierda) saca de los registros las double words calculadas en el paso anterior con los $u$ y $v$ y chequea de a una las restricciones del cálculo para $u$ y $v$, para luego colocar en el caso de que las cumplan, en el acumulador del ciclo, el valor del pixel que tenga en la posición en la imagen fuente (el cual es leído desde esta) y si no, un 0.

Terminado el ciclo de almacenamiento, se agregan los valores de estos 4 bytes al acumulador principal creado al comienzo del ciclo, utilizando \texttt{PSLLDQ} para shiftearlo a la izquierda 4 bytes y agregando los valores mediante un por y se vuelve a ciclar realizando el mismo proceso hasta llenar el acumulador con los valores correspondientes de los 16 bytes o 16 píxeles contiguos.


\subsubsection{Escritura de los píxeles en la imagen destino y fin del ciclo principal}

En el final de cada iteración del ciclo principal, se chequea la posición de la imagen en la que se está (por si ya se ha terminado de recorrerla o se es necesario correr para atrás para agarrar 16 píxeles exactos) y si se ha llenado el acumulador. Si el acumulador se llena (habiendo pasado 4 ciclos) se escriben los 16 bytes a la imagen destino. Este mecanismo en el cuál se van almacenando los píxeles tiene como sentido economizar la escritura de datos.


\subsubsection{Comparación con la implementación C}

Cómo el algoritmo no es muy paralelizable, ambas implementaciones comparten varias características pero difieren en algunos casos, a continuación se detalla.

Operaciones realizadas en la implementación en C para 16 bytes:
\begin{itemize}
\item Se realizan $2 * 16 = 32$ accesos a memoria, 16 para la lectura de un pixel determinado en la imagen fuente y otros 16 para la escritura en la imagen destino.
\item Se realiza el cálculo relacionado a $u$ y $v$ de manera secuencial.
\item Se realiza el chequeo de las restricciones para $u$ y $v$ de manera secuencial.
\end{itemize}

Operaciones realizadas en la implementación en ASM para 16 bytes:
\begin{itemize}
\item Se realizan 16 accesos a memoria para la lectura de un pixel determinado en la imagen fuente y 1 acceso utilizando \asm{MOVDQU} para la escritura en la imagen destino.
\item Se realiza el cálculo relacionado a $u$ y $v$ utilizando instrucciones SIMD pudiendo realizar los cálculos con 2 datos simultáneos realizando en cada ciclo 4 cálculos de $u$ y $v$ y en cada ciclo principal el cálculo de los 16 bytes.
\item Se realiza el chequeo de las restricciones para $u$ y $v$ de manera secuencial.
\end{itemize}


\subsubsection{Rendimiento}

Observamos las siguientes cantidades de ciclos y ticks de reloj al realizar 100 iteraciones de ambas implementaciones con una imagen cuadrada de lado 512.
\begin{center}
    \begin{tabular}{|l|l|l|l|}
        \hline
        Medición & Implementación C & Implementación assembler & Relación \\
        \hline
        Ticks    & 1766855008      & 727560304               & $41.17\%$ \\
        Ciclos   & 17668550       & 7275603                & $41.17\%$ \\
        \hline
    \end{tabular}
\end{center}

Es notable la diferencia en velocidad con los otros filtros. Este resultado pone en evidencia que al poder paralelizarse en menor manera que otros filtros, la velocidad y los ciclos aumentan considerablemente.

Cabe destacar que la diferencia de velocidad obtenida en comparación con la implementación en C no sólo proviene de el uso de instrucciones SIMD para el cálculo con doubles o el acceso a memoria usando \asm{MODQU}, si no que también influyen el uso continuo de registros en vez de accesos memoria para almacenar datos (algo que C hace) y la posibilidad que otorga el lenguaje ensamblador de ordenar el código para poder optimizar al máximo los algoritmos.

%%%%%%%%%%%%%%%%%%%%%%%%%%%%%%%%%%%%%%%%%%%%%%%%%%%%%%%%%%%%%%%%%%%%%%%%%%%%%%%
%% Resultados                                                                %%
%%%%%%%%%%%%%%%%%%%%%%%%%%%%%%%%%%%%%%%%%%%%%%%%%%%%%%%%%%%%%%%%%%%%%%%%%%%%%%%


\section{Resultados}

Pendiente.


%%%%%%%%%%%%%%%%%%%%%%%%%%%%%%%%%%%%%%%%%%%%%%%%%%%%%%%%%%%%%%%%%%%%%%%%%%%%%%%
%% Conclusión                                                                %%
%%%%%%%%%%%%%%%%%%%%%%%%%%%%%%%%%%%%%%%%%%%%%%%%%%%%%%%%%%%%%%%%%%%%%%%%%%%%%%%


\section{Conclusión}

Podemos concluir, mediante los resultados obtenidos para cada filtro, que el uso de instrucciones SIMD mejora considerablemente el uso de ciclos y el tiempo, poniendo en evidencia el potencial de estas instrucciones.

También pudimos observar desensamblando los binarios, el código ensamblador generador por gcc, el cuál puso en evidencia la gran cantidad de accesos a memoria que utiliza el código generado por C debido al uso de variables locales. Estos accesos que fueron evitados mediante el uso de registros en las versiones ASM, decrementaron en cada filtro el número de ciclos utilizados ya que la velocidad de operación sobre los registros es mucho más rápida que la de memoria.

Por último, pudimos observar en la realización de las implementaciones en ASM que no todos los algoritmos son buenos candidatos para la parelización. Algoritmos como umbralizar o colorizar fueron buenos candidatos para la parelización pero otros como rotar, no tienen la posibilidad de aprovechar al máximo la paralelización.


\end{document}
